\chapter{Geometry from Linear Structure}

In Chapter I, vector spaces were introduced as purely algebraic objects.
A vector was an abstract element of a set equipped with addition and scalar multiplication,
and linear transformations were functions preserving this structure.
No notion of length, angle, distance, or orthogonality was assumed or required.

\medskip

Algebra alone allows us to decide whether vectors are independent,
whether a set spans a space,
or whether a linear transformation is injective or surjective.
However, it does not allow us to ask geometric questions.
In many situations, it is not enough to know \emph{which} vectors are possible;
we also need to know how \emph{large} they are, how \emph{close} they are to one
another, or how \emph{aligned} they are.
These notions are expressed through length, distance, and angle.
Without additional structure, the language of vector spaces does not distinguish
between ``larger'' and ``smaller,'' nor does it allow one to speak of
perpendicularity or of the ``best approximation'' of one vector by others
(as in projections and least squares).
There is no intrinsic way, within a bare vector space, to compare the size of
two vectors or to determine whether they point in perpendicular directions.

\medskip

Geometry begins only when additional structure is imposed.

\medskip

It is useful now to recall how geometric space is commonly represented in
coordinates.
Given sets $A_1,\dots,A_n$, their \emph{Cartesian product}
(also known as \emph{Product Set}) is
\[
A_1\times\cdots\times A_n = \{(x_1,\dots,x_n)\mid x_i\in A_i\}.
\]
In particular, the coordinate space $\mathbb{R}^n$ is the Cartesian product of
$n$ copies of $\mathbb{R}$.
This construction by itself is not yet geometry: it provides a convenient
\emph{model} in which ``points'' are represented by tuples of real numbers.
A geometry emerges only after we specify an additional rule that allows us to
measure and compare such tuples.

\medskip

A central theme of this document is that geometry is not a replacement for
linear structure, but an enrichment of it.
The vector spaces introduced earlier remain unchanged;
what changes is that they are now equipped with a rule that allows vectors to
interact in a way that produces scalar quantities.

\medskip

Informally, this rule takes two vectors and returns a scalar.
From this single operation, notions such as length, angle, and distance emerge
naturally.

\medskip

To make this interaction precise, we introduce the notion of a bilinear map.

\medskip

Let $E$ and $G$ be vector spaces over the same field $K$.
A mapping
\[
B : E \times G \longrightarrow K
\]
is called \emph{bilinear} if it is linear in each argument separately.
That is, for fixed $y \in G$, the map
\[
x \longmapsto B(x,y)
\]
is a linear functional on $E$, and for fixed $x \in E$, the map
\[
y \longmapsto B(x,y)
\]
is a linear functional on $G$.

\medskip

Explicitly, for all $x_1,x_2 \in E$, $y_1,y_2 \in G$, and all scalars
$\alpha \in K$,
\[
B(x_1 + x_2, y) = B(x_1, y) + B(x_2, y), \qquad
B(\alpha x, y) = \alpha B(x, y),
\]
\[
B(x, y_1 + y_2) = B(x, y_1) + B(x, y_2), \qquad
B(x, \alpha y) = \alpha B(x, y).
\]

\medskip

When $E = G$, such a mapping is called a \emph{bilinear form} on $E$.

\medskip

At this level, no geometric interpretation is assumed.
A bilinear form is merely a rule that assigns a scalar to an ordered pair of
vectors, subject to linearity in each argument.
It need not be symmetric, positive, or related to any notion of length or angle.

\medskip

Nevertheless, bilinear forms play a fundamental structural role.
They provide a systematic way to associate vectors with linear functionals
and to construct scalar quantities from pairs of vectors.

\medskip

Only after additional conditions are imposed on a bilinear form do familiar
geometric notions emerge.
These special cases will be introduced shortly.

% CORRECTION 3: The restriction to R is now introduced here, immediately after
% bilinear forms and before any mention of positivity or inner products.
% Previously it appeared much later, after the inner product had already been
% implicitly used.

\medskip

Up to this point, the scalar field $K$ has been arbitrary.
Pure linear structure does not require the ability to compare scalars.
Geometry, however, depends on positivity: expressions such as
$\langle x,x\rangle > 0$ and inequalities of the form $r>0$
presuppose an order relation.
For this reason, we now restrict attention to the real number system
\[
\Big\langle \mathbb{R}, +, \cdot, \le \Big\rangle,
\]
viewed as an ordered field.
It is precisely the order on $\mathbb{R}$ that allows lengths, distances, and
orthogonality to acquire geometric meaning.

\medskip

We now specialize the previously introduced bilinear forms.

\medskip

Let $E$ be a vector space over $\mathbb{R}$.
A \emph{bilinear form} on $E$ is a function
\[
B : E \times E \longrightarrow \mathbb{R}
\]
that is linear in each argument separately.
That is, for all $x,x',y,y' \in E$ and all $\lambda \in \mathbb{R}$,
\[
B(x+x',y)=B(x,y)+B(x',y), \qquad
B(\lambda x,y)=\lambda B(x,y),
\]
and similarly in the second argument.

\medskip

Bilinear forms arise naturally from linear transformations.
If
\[
\psi : E \longrightarrow E^*
\]
is a linear mapping, then the formula
\[
B(x,y) = \psi(x)(y)
\]
defines a bilinear form on $E$.
Conversely, every bilinear form determines such a mapping.
Therefore, bilinear forms encode a controlled way in which vectors may interact
to produce scalars.

\medskip

At this level, no geometry is present.
A bilinear form may be degenerate, asymmetric, or indefinite.
It allows vectors to interact, but it does not yet measure anything.

\medskip

Geometry appears when the bilinear form satisfies additional constraints.

% CORRECTION 1 & 2: The inner product is now defined before any coordinate
% expressions involving it. The section on the transpose, which previously
% interrupted the path from bilinear forms to the inner product, has been
% relocated to after the inner product is established. The coordinate expression
% of the inner product follows the definition, not precedes it.

\medskip

An \emph{inner product} on a vector space $E$ over $\mathbb{R}$
is a bilinear form
\[
\langle \cdot,\cdot \rangle : E \times E \longrightarrow \mathbb{R}
\]
that is symmetric, positive, and non-degenerate.
These properties ensure that
\[
\langle x,x \rangle > 0 \quad \text{for all } x \neq 0.
\]

\medskip

Once such a rule is fixed, vectors acquire measurable properties.
Two vectors may be declared orthogonal.
A vector may be assigned a length.
The distance between two points may be defined in terms of the vectors
connecting them.
None of these concepts exist prior to the introduction of this additional
structure.

\medskip

As in Chapter I, coordinates play no fundamental role.
Geometric quantities must be independent of any particular reference system.
A formula that changes when the basis is changed describes a representation,
not a geometric object.

\medskip

Only after a basis is chosen do familiar coordinate expressions appear.
Sums of squares, dot products, and quadratic equations arise as coordinate-level
manifestations of the underlying structure, not as definitions of it.

\medskip

\textit{Coordinate expression of the inner product.}

Let $(E,\langle\cdot,\cdot\rangle)$ be a finite-dimensional inner product space
and let $\{e_1,\dots,e_n\}$ be an orthonormal basis.
Every vector $x\in E$ can be written uniquely as
\[
x=\sum_{i=1}^n x_i e_i,
\qquad
y=\sum_{i=1}^n y_i e_i.
\]

Using bilinearity and the orthonormality condition
$\langle e_i,e_j\rangle=\delta_{ij}$, we obtain
\[
\langle x,y\rangle
=
\sum_{i=1}^n x_i y_i.
\]

Thus, in an orthonormal basis the abstract inner product is represented by the
familiar \emph{dot product} of coordinate vectors in $\mathbb{R}^n$.

\medskip

This observation explains the role of matrix multiplication.
If a linear transformation $\psi:E\to G$ is represented by a matrix
$A=(a_{ij})$ and a vector $x$ has coordinate vector $x=(x_1,\dots,x_n)^T$,
then the $i$-th component of the product $Ax$ is
\[
(Ax)_i=\sum_{j=1}^n a_{ij}x_j.
\]

Each component is therefore the dot product of the $i$-th row of $A$
with the vector $x$.
Matrix multiplication can thus be understood as the systematic evaluation of
inner products between rows of one matrix and columns of another.

\medskip

In this sense, the familiar algebraic rule for multiplying matrices is not an
arbitrary computational convention: it admits a natural geometric interpretation
as the coordinate expression of interactions measured by an inner product.

\medskip

This point of view explains why objects such as circles, spheres, and ellipsoids
can be described by algebraic equations, yet retain a geometric meaning that is
invariant under changes of coordinates.
The equations depend on the chosen basis; the geometry does not.

\medskip

The goal is therefore to show how geometry emerges from linear structure once a
suitable interaction between vectors is specified.

\medskip

% CORRECTION 2 (continued): The section on the transpose and dual transformation
% is now placed here, after the inner product is defined. This preserves the
% algebraic derivation of the transpose as intrinsic (no inner product needed)
% while allowing the remark about duality becoming internal to appear naturally
% in the context of an established inner product.

Although geometry has begun to emerge through bilinear interaction, the
underlying algebraic mechanism remains the linear transformation.
To understand how geometric structure interacts with linear maps, we examine a
construction naturally associated with every linear transformation.

\medskip

The central object of this work is the linear transformation. Before introducing
geometric notions such as length, angle, and distance, we make explicit a
fundamental construction associated with every linear transformation: its
transpose.

\medskip

Let $E$ and $G$ be vector spaces over the same field $K$, and let
\[
\psi : E \longrightarrow G
\]
be a linear transformation. Recall that the dual spaces $E^*$ and $G^*$ consist
of all linear functionals on $E$ and $G$, respectively.

\medskip

For each $g^* \in G^*$, the composition
\[
g^* \circ \psi : E \longrightarrow K
\]
is a linear functional on $E$. This defines a mapping
\[
\psi^* : G^* \longrightarrow E^*,
\qquad
\psi^*(g^*) = g^* \circ \psi.
\]

\medskip

The mapping $\psi^*$ is linear and is called the \emph{transpose}
(or dual transformation) of $\psi$. The direction of the arrow is reversed:
\[
E \xrightarrow{\;\psi\;} G
\quad\Longrightarrow\quad
G^* \xrightarrow{\;\psi^*\;} E^*.
\]

\medskip

This construction is intrinsic: it depends only on the linear transformation
$\psi$ and does not require any choice of basis, coordinates, or inner product.

\medskip

When bases are chosen for $E$ and $G$, and the corresponding dual bases are
chosen for $E^*$ and $G^*$, the linear transformation $\psi$ is represented by
a matrix $A$, and the dual transformation $\psi^*$ is represented by the
transposed matrix $A^T$.
Thus, the familiar matrix transpose is the coordinate expression of the intrinsic
dual transformation $\psi^*$.

\medskip

The inner product does more than assign lengths.
It establishes a canonical identification between the vector space and its dual.

\medskip

For each vector $x \in E$, define the mapping
\[
\Phi(x) : E \longrightarrow \mathbb{R},
\qquad
\Phi(x)(y) = \langle x,y \rangle.
\]

By bilinearity, $\Phi(x)$ is a linear functional on $E$.
Thus, we obtain a mapping
\[
\Phi : E \longrightarrow E^*.
\]

\medskip

If the inner product is non-degenerate, this mapping is injective.
In finite dimensions, since $\dim E = \dim E^*$, it is therefore an isomorphism.

\medskip

Hence, once an inner product is fixed, every vector determines a unique linear
functional, and every linear functional arises from a unique vector.
The distinction between $E$ and $E^*$ disappears at the structural level.

\medskip

It is precisely this identification that allows the transpose of a linear
transformation to be regarded as an operator on $E$ itself.
When geometry is introduced, duality becomes internal.

\medskip

Once an inner product is fixed, vectors acquire length.
The norm of a vector $x \in E$ is defined by
\[
\|x\| = \sqrt{\langle x,x \rangle}.
\]

\medskip

So, the vector space $E$ introduced in Chapter I is now equipped with additional
geometric structure:
\[
\Big\langle
\langle \mathbb{R},+,\cdot\rangle,\;
\langle E,+\rangle,\;
\circ,\;
\langle\cdot,\cdot\rangle
\Big\rangle.
\]

The inner product is not required for linearity itself, but it allows the
introduction of geometric notions such as length, angle, orthogonality, and
adjoint transformations.

\medskip

\textit{Algebraic expansion of the norm.}

For any vectors $x,y \in E$, the square of the norm of their sum satisfies
\[
\|x+y\|^2
=
\langle x+y, x+y \rangle.
\]
By bilinearity and symmetry of the inner product,
\[
\langle x+y, x+y \rangle
=
\langle x,x \rangle
+
\langle x,y \rangle
+
\langle y,x \rangle
+
\langle y,y \rangle
=
\|x\|^2
+
2\langle x,y \rangle
+
\|y\|^2.
\]

Thus,
\[
\|x+y\|^2
=
\|x\|^2
+
\|y\|^2
+
2\langle x,y \rangle.
\]

This identity holds in every inner product space.
It expresses how the interaction term $\langle x,y \rangle$ measures the
deviation from pure additivity of squared lengths.

\medskip

The inner product also satisfies the inequality
\[
|\langle x,y \rangle|
\le
\|x\|\,\|y\|,
\]
known as the \emph{Cauchy--Schwarz inequality}.
As a consequence,
\[
\|x+y\|^2
\le
\|x\|^2
+
\|y\|^2
+
2\|x\|\,\|y\|
=
(\|x\|+\|y\|)^2.
\]

Taking square roots yields the \emph{triangle inequality}
\[
\|x+y\|
\le
\|x\|
+
\|y\|.
\]

\medskip

The \textit{Pythagorean identity} arises as the special case
$\langle x,y \rangle = 0$,
in which the interaction term vanishes and
\[
\|x+y\|^2
=
\|x\|^2
+
\|y\|^2.
\]

\medskip

Orthogonality also becomes meaningful.
Two vectors $x,y \in E$ are said to be orthogonal if
\[
\langle x,y \rangle = 0.
\]

\medskip

The order of the vectors is irrelevant.
Symmetry of the inner product implies
\[
\langle x,y \rangle = \langle y,x \rangle,
\]
so orthogonality is a mutual relation.

\medskip

The zero vector is orthogonal to every vector in $E$.
Indeed, for any $v \in E$,
\[
\langle 0,v \rangle
=
\langle 0\cdot v, v \rangle
=
0\langle v,v \rangle
=
0.
\]
Moreover, the zero vector is the only vector with this property.

\medskip

Thus, orthogonality is precisely the condition under which the interaction term
vanishes.
In that case,
\[
\|x+y\|^2 = \|x\|^2 + \|y\|^2.
\]
Squared lengths become additive exactly when the vectors do not interact.

\medskip

% CORRECTION 4: The interlude on eigenvectors has been moved to here, after
% the norm, orthogonality, and the adjoint have been established. Previously
% it appeared between the Pythagorean identity and the section on distance,
% interrupting the chain norm -> distance -> topology at its most critical point.

\textit{Interlude: Invariant Directions}

Up to this point, linear transformations have been treated algebraically:
a map $\psi : E \to E$ sends vectors to vectors while preserving linear
structure.
No geometric meaning has yet been attached to this action.

\medskip

Once an inner product is introduced, however, the action of a linear
transformation can be examined geometrically.
Given a vector $x \in E$, the image $\psi(x)$ may differ from $x$ in two
independent ways: its length may change, and its direction may change.

\medskip

For a generic vector, both effects occur simultaneously.
Most vectors are stretched and rotated.

\medskip

A remarkable phenomenon occurs when there exists a nonzero vector $x \in E$
such that
\[
\psi(x) = \lambda x
\]
for some scalar $\lambda$.
In this case, the transformation acts on $x$ by pure scaling.
The direction of $x$ is preserved.

\medskip

Such a vector is called an \emph{eigenvector} of $\psi$, and the scalar
$\lambda$ is its corresponding \emph{eigenvalue}.
Eigenvectors identify directions in the space that are geometrically invariant
under the transformation.

\medskip

This invariance is an intrinsic property of the linear transformation itself.
Changing the basis may alter the representation of $\psi$, but it cannot
destroy or create eigenvectors.

\medskip

Under the identification $E \cong E^*$ induced by the inner product, the
transpose $\psi^*$ corresponds to a linear operator on $E$, called the
\emph{adjoint} of $\psi$.

\medskip

When a linear transformation arises from a symmetric bilinear form (or
equivalently, from a self-adjoint operator), its eigenvectors corresponding
to distinct eigenvalues are orthogonal.
These directions define a preferred coordinate system dictated by the geometry,
not by arbitrary choice.

\medskip

It is in this sense that eigenvectors serve as the principal axes of the
ellipsoids associated with quadratic forms.
Diagonalization is the act of aligning coordinates with the invariant directions
already present in the structure.

\medskip

% CORRECTION 4 (continued): With eigenvectors established here, the section on
% distance and topology now follows without interruption.

\textit{Distance and the emergence of geometry.}

Once an inner product is available on a vector space $E$, it determines a norm,
and from the norm we obtain a notion of distance between vectors,
\[
d(x,y)=\|x-y\|.
\]

At this point, no new structure is introduced: distance is merely a
reformulation of algebraic data already present.

\medskip

This distance allows us to describe proximity. For a vector $x \in E$ and a
real number $r > 0$, we consider the set
\[
B(x,r) = \{\, y \in E \mid d(x,y) < r \,\}.
\]
Such a set is called an \textit{open ball} centered at $x$ with radius $r$.
Geometrically, it consists of all vectors that lie closer to $x$ than the
prescribed threshold.

\medskip

The strict inequality is essential. It guarantees that every point inside the
ball admits further nearby points, so that no boundary is included.
This ``breathing room'' is what ultimately allows us to speak of continuity and
deformation.

\medskip

The collection of all open balls determines which subsets of $E$ should be
called open: a set is open if, around each of its points, it contains some open
ball.
In this way, distance generates a notion of openness without any additional
axioms.

% CORRECTION 5: A transitional sentence now prepares the reader for the
% axiomatic definition of a topological space, connecting it to the geometric
% intuition built just above rather than introducing the axioms abruptly.

\medskip

This notion of openness can be captured by a small set of axioms that
preserve its essential properties while allowing the concept to be extended
far beyond the metric setting.
The resulting structure is what is known as a \textit{topology} on $E$.

\medskip

A \textit{topological space} is an ordered pair $(X, \tau)$, where $X$ is a
set and $\tau$ is a collection of subsets of $X$.
For $\tau$ to be a \textit{topology}, it must satisfy the following axioms:

\begin{enumerate}
    \item $\emptyset \in \tau$ and $X \in \tau$.
    \item The union of any collection of sets in $\tau$ is also in $\tau$.
    \item The intersection of any \textit{finite} number of sets in $\tau$ is
    also in $\tau$.
\end{enumerate}

The elements of $\tau$ are called \textit{open sets}.

\medskip

Thus, topology does not appear here as an independent theory, but as a
structural consequence of the metric induced by the inner product.
Geometry gives rise to distance, distance gives rise to neighborhoods, and
neighborhoods give rise to topological structure.

\medskip

\textit{Geometric interpretation.}

In $\mathbb{R}^n$ equipped with its standard inner product, open balls
correspond to the familiar spheres of Euclidean geometry.
More general norms deform these balls into ellipsoidal shapes, reflecting
anisotropic scaling along different directions of the space.
In all cases, each point
\[
x = (x_1,\dots,x_n)
\]
is a vector of the ambient vector space, and the geometric objects considered
are subsets of that space defined by distance relations among its vectors.

\medskip

Once objects have been represented as vectors in this way, the geometric
structures introduced above---norms, distances, balls, and ellipsoids---become
tools for organizing and comparing them.

\medskip

Their relative positions, measured by a chosen metric, encode similarity.
This metric need not be Euclidean: in practice it may be designed to reflect
domain-specific invariances or even learned from data.
Neighborhoods of meaning are then described by open balls, and learning can be
interpreted as the discovery of a geometric and topological organization
underlying the data.

\medskip

Once a norm has been introduced, certain geometric sets arise as direct
consequences of the structure and require no additional assumptions.

\medskip

Let $(E,\langle\cdot,\cdot\rangle)$ be an inner product space and let
$\|x\|=\sqrt{\langle x,x\rangle}$ be the associated norm.
For a fixed vector $x_0\in E$ and a scalar $r>0$, the \emph{closed ball}
of radius $r$ centered at $x_0$ is defined by
\[
B_r(x_0)=\{\,x\in E \mid \|x-x_0\|\le r\,\}.
\]

The corresponding \emph{sphere} is the boundary of this set,
\[
S_r(x_0)=\{\,x\in E \mid \|x-x_0\|= r\,\}.
\]

\textit{Remark (Coordinate description in $\mathbb{R}^n$):}
When $E=\mathbb{R}^n$ is equipped with its standard inner product and we use
the standard coordinates $x=(x_1,\dots,x_n)$, the same intrinsic definitions
take the familiar ``sum of squares'' form.
For a ball centered at the origin,
\[
B_r(0)=\left\{x\in\mathbb{R}^n\;\middle|\;x_1^2+\cdots+x_n^2\le r^2\right\},
\]
and its boundary sphere is
\[
S_r(0)=\left\{x\in\mathbb{R}^n\;\middle|\;x_1^2+\cdots+x_n^2=r^2\right\}.
\]
This is simply the coordinate expression of
$\|x\|=\sqrt{x_1^2+\cdots+x_n^2}$.

\medskip

These definitions are intrinsic and do not depend on any choice of coordinates
or basis.

\medskip

\textit{Remark (Symmetry):}
In the Euclidean case $E=\mathbb{R}^n$ with its standard inner product, the
ball is invariant under orthogonal transformations.
If $Q\in O(n)$, then $\|Qx\|=\|x\|$, and hence
\[
Q\bigl(B_r(0)\bigr)=B_r(0).
\]
More generally, any isometry of an inner product space carries balls to balls.

\medskip

\textit{Remark (Slicing in $\mathbb{R}^n$):}
In $\mathbb{R}^n$, fixing the last coordinate $x_n=t$ shows that the
cross-section of a radius-$r$ ball by the hyperplane $x_n=t$ is an
$(n-1)$-dimensional ball of radius $\sqrt{r^2-t^2}$.
This observation is the geometric content behind the common
``integrate by slices'' approach to $n$-dimensional volume.

\medskip

\textit{Remark (Volume of Euclidean balls):}
In $\mathbb{R}^n$, the $n$-dimensional volume of a Euclidean ball scales as a
constant times $r^n$:
\[
\operatorname{Vol}_n\bigl(B_r(0)\bigr)=\omega_n\,r^n,
\]
where $\omega_n=\operatorname{Vol}_n\bigl(B_1(0)\bigr)$ is the volume of the
unit ball. A closed form is
\[
\omega_n=\frac{\pi^{n/2}}{\Gamma\!\left(\frac{n}{2}+1\right)}.
\]
In particular, $\omega_1=2$, $\omega_2=\pi$, and $\omega_3=\frac{4\pi}{3}$.

\medskip

\textit{Remark (A high-dimensional effect):}
Although $B_1(0)\subset \mathbb{R}^n$ always contains many points, its
$n$-dimensional volume satisfies $\omega_n\to 0$ as $n\to\infty$.
This illustrates a basic phenomenon of high-dimensional geometry: sets that
look large in low dimensions can occupy a vanishing fraction of ambient volume
as the dimension increases.

\medskip

More generally, let $B:E\times E\to \mathbb{R}$ be a symmetric
positive-definite bilinear form. The set
\[
\mathcal{E}=\{\,x\in E \mid B(x-x_0,x-x_0)\le r^2\,\}
\]
is called an \emph{ellipsoid}.
Thus, an ellipsoid is the unit ball associated with a bilinear form, or
equivalently, with a modified inner product on $E$.

\medskip

When the bilinear form $B$ is symmetric and positive-definite, there exists a
unique self-adjoint linear operator $A:E\to E$ such that
\[
B(x,y)=\langle Ax,y\rangle
\]
for all $x,y\in E$.

\medskip

The operator $A$ encodes the geometry of the ellipsoid.
Since $A$ is self-adjoint on a finite-dimensional inner product space, it
admits an orthonormal basis of eigenvectors.
If $\{e_1,\dots,e_n\}$ is such a basis and $A e_i=\lambda_i e_i$ with
$\lambda_i>0$, then the defining inequality
\[
B(x-x_0,x-x_0)\le r^2
\]
becomes
\[
\sum_{i=1}^n \lambda_i\, \xi_i^2 \le r^2,
\]
where $\xi_i=\langle x-x_0,e_i\rangle$ are the coordinates in this
eigenbasis.

% CORRECTION 6: The observation that diagonalization "reveals invariant
% directions already present" appeared twice in close succession. The first
% occurrence (here, in the ellipsoid derivation) is retained as it concludes
% the algebraic argument. The redundant restatement that appeared a few lines
% later, just before the coordinate formula, has been removed.

\medskip

Thus, the principal axes of the ellipsoid are precisely the eigenvectors of
$A$, and the eigenvalues determine the scaling along those directions.
Diagonalization does not create the axes; it merely reveals the invariant
directions already present in the structure.

\medskip

Accordingly, when a basis is chosen in which the bilinear form is diagonal,
this invariant definition reduces to a familiar coordinate expression.
In $\mathbb{R}^n$, one obtains
\[
\mathcal{E}
=
\left\{
(x_1,\dots,x_n)\in\mathbb{R}^n
\;\middle|\;
\sum_{i=1}^n \frac{(x_i-a_i)^2}{h_i^2} \le r^2
\right\}.
\]

Written componentwise, this inequality reads
\[
\left(\frac{x_1-a_1}{h_1}\right)^2
+
\left(\frac{x_2-a_2}{h_2}\right)^2
+
\cdots
+
\left(\frac{x_n-a_n}{h_n}\right)^2
\le r^2.
\]

The equation depends on the chosen coordinates; the ellipsoid itself does not.

\medskip

\emph{Remark (Degenerate cases).}
The geometric nature of the set defined above depends on the coefficients
$h_1,\dots,h_n$.
If $h_1=\cdots=h_n$, the ellipsoid reduces to an $n$-dimensional sphere.
If one or more of the parameters $h_i$ vanish, the defining quadratic form
becomes degenerate and the ellipsoid collapses into a lower-dimensional
object: a disk, a line segment, or, in the extreme case, a single point.

\medskip

From the structural point of view, these degenerations correspond to a loss of
rank in the associated quadratic or bilinear form.
Thus, changes in the algebraic structure translate directly into changes in
geometric dimension.

\medskip

\emph{Remark.}
Throughout this discussion, the ambient space $E=\mathbb{R}^n$ is a vector
space in the sense of Chapter I.
Each element
\[
x=(x_1,\dots,x_n)
\]
is therefore a vector of $E$.

\medskip

However, the sets defined by these expressions are \emph{subsets} of $E$
selected by quadratic constraints.
They are not vector spaces themselves, since they are not closed under
addition or scalar multiplication.

\bigskip

\textit{Conceptual summary.}
Vector spaces provide the stage.
Linear transformations describe allowable changes of state.
Geometry enters only when vectors are permitted to interact through a
structured bilinear rule that produces scalars.
From this interaction, lengths, angles, distances, and geometric shapes emerge.
The transition from algebra to geometry is therefore not a leap, but a
controlled enrichment of structure.
The book \cite{iribarren2} is a wonderful reference on the topology of metric
spaces.

\section*{Notes and Remarks}

\textit{Outlook.}

In contemporary applications such as representation learning, elements of a
set (words, images, or signals) are mapped into high-dimensional vector spaces,
where geometric relations encode similarity and structural information.
These constructions provide concrete realizations of abstract linear-algebraic
ideas.

\medskip

Modern computational systems offer another illustrative example.
Platforms devoted to \emph{observability}---that is, the systematic measurement
of the internal state of a running system---collect large families of numerical
quantities such as latency, throughput, memory consumption, error rates, and
resource utilization.

\medskip

At a fixed instant of time, these measurements may be assembled into a vector
\[
x(t) = (x_1(t),\dots,x_n(t)) \in \mathbb{R}^n,
\]
whose coordinates represent observable characteristics of the system.
As time evolves, the system determines a trajectory in this ambient linear
space, so that monitoring and analysis may be interpreted geometrically as
the study of paths in a finite-dimensional vector space.

\medskip

Operational considerations typically impose bounds on each observable,
expressed by inequalities of the form
\[
a_i \le x_i \le b_i,
\]
reflecting physical limitations or acceptable operating ranges.
Geometrically, these constraints define a Cartesian product of intervals,
namely an $n$-dimensional hyperrectangle (or, after normalization, a
hypercube).
Normal system behavior may therefore be viewed as confined to a bounded region
of the ambient space, while anomalous behavior corresponds to trajectories
departing from this region.

\medskip

It should be emphasized that this construction is an idealized mathematical
model: in practical systems the collection of observables may vary over time,
so the dimension $n$ is best regarded as fixed only locally or within a chosen
representation.
Nevertheless, the linear-space framework provides a natural language for
describing states, constraints, and evolution.

\medskip

In this way, observability systems illustrate a recurring theme of linear
algebra: a flat ambient space serves as a universal stage upon which meaningful
phenomena are characterized by geometrically defined subsets.
Norm constraints give rise to spheres, anisotropic scalings produce ellipsoids,
and coordinate-wise bounds define hyperrectangles.
Despite their differing geometric appearances, all arise as structured subsets
embedded within the same linear ambient space.

\bigskip
\hrule
\medskip

% CORRECTION 7: The historical remark on Lambert has been moved to immediately
% after the Interlude on Invariant Directions in the Notes section, where its
% connection to eigenvectors and angle preservation is contextually immediate.
% Previously it appeared later in the notes, separated from the discussion that
% motivated it.

\textit{Historical remark (Angles and invariance).}

Before eigenvectors were understood as invariant directions, their detection
was often indirect.
In two dimensions, determining whether a transformation preserves a direction
reduces to detecting whether it preserves an angle.

\medskip

This problem led naturally to the computation of arctangents.
In the eighteenth century, J.H.~Lambert introduced the continued fraction
\[
\arctan(z)
=
\cfrac{z}{1 + \cfrac{z^2}{3 + \cfrac{4z^2}{5 + \cfrac{9z^2}{7 + \cdots}}}},
\]
as part of his work on orbital mechanics and rotational phenomena.

\medskip

From the modern point of view, this analytic effort reflects a geometric
question: does a linear transformation rotate a direction, or does it merely
scale it?
Eigenvectors provide the structural answer to this question, rendering angle
computation auxiliary rather than fundamental.

\medskip
\hrule
\bigskip

\textit{Remark (Quadratic forms and the chi-square statistic).}
Expressions of the form defining ellipsoids appear naturally in contexts that
are not, at first sight, geometric.
A notable example arises in the theory of random number testing, where one
encounters the quantity
\[
V
=
\frac{(Y_1 - np_1)^2}{np_1}
+
\frac{(Y_2 - np_2)^2}{np_2}
+
\cdots
+
\frac{(Y_n - np_n)^2}{np_n},
\]
known in statistics as the \emph{chi-square statistic} associated with the
observed quantities $Y_1,\dots,Y_n$ and reference values $np_1,\dots,np_n$.

\medskip

From the present point of view, this expression is simply a quadratic form on
$\mathbb{R}^n$.
It measures the squared length of the deviation vector
\[
(Y_1 - np_1,\dots,Y_n - np_n)
\]
with respect to a symmetric positive-definite bilinear form whose matrix, in
the standard basis, is diagonal with entries
$(np_1)^{-1},\dots,(np_n)^{-1}$.

\medskip

Consequently, the sets defined by inequalities of the form
\[
V \le c
\]
are ellipsoids centered at the point $(np_1,\dots,np_n)$.
The denominators $np_i$ determine the anisotropic scaling along the coordinate
directions, exactly as in the general ellipsoid
\[
\sum_{i=1}^n \frac{(x_i-a_i)^2}{h_i^2} \le r^2.
\]

\medskip

Thus, before any probabilistic interpretation is imposed, the chi-square
expression is already a geometric object: it is the squared norm induced by a
particular inner product on $\mathbb{R}^n$.
Statistical theory makes use of this geometry, but does not create it.

\medskip

\textit{Remark (The ambient vector space and constrained regions).}
The vector space $E$ itself is the \emph{ambient space} in which all vectors
live.
In coordinates, when a basis is chosen, this space appears as $\mathbb{R}^n$,
a flat linear environment in which vectors may be freely added and scaled.

\medskip

Geometric objects such as balls, spheres, and ellipsoids do not replace this
ambient space; rather, they are subsets of it defined by constraints on the
norm or on a quadratic form.
For example, the sphere
\[
S_r(x_0)=\{x\in E \mid \|x-x_0\|=r\}
\]
selects vectors whose distance from $x_0$ is fixed, while an ellipsoid selects
vectors whose coordinates satisfy an anisotropic quadratic inequality.

\medskip

In many modern applications (AI Deep Learning for instance), algorithms operate
on vectors in a high-dimensional ambient space while imposing such constraints
during computation.
Normalization steps may restrict vectors to spheres, regularization procedures
may confine parameters to ellipsoidal regions, and optimization methods search
for solutions within neighborhoods described by balls.
The ambient vector space provides the linear universe in which vectors exist,
while these geometric subsets describe regions selected by the rules or
constraints of the problem.